% =====================================================================
%  Chapter 25 opener -- "Chapter at a Glance", three reading paths + the
%  ladder staircase, THE running example (the shared 8-node rooted forest R),
%  the who-owns-what comparison against Ch.8 (splay intro), Ch.22 (HLD),
%  Ch.16 (LCA), Ch.5 (Euler tour) and Ch.24 (centroid), the recognition
%  guide, a constraint->technique table, and notation.
%  Exports fig:c23-tree, fig:c23-taxonomy.
%  No \term{splay tree} here -- it is indexed by Ch.8 (\chref{c06}); \term the
%  NEW objects "link-cut tree", "preferred path", "represented tree" only.
% =====================================================================

\begin{chaptertheorybox}{Chapter at a Glance}
Chapter~8 showed you a \emph{splay tree}: a binary search tree that rotates the
node you just touched up to the root, so frequently-used keys float to the top and
a long run of accesses costs \(O(\log n)\) each \emph{amortised} (\secref{c06-splay}).
It closed with an explicit IOU -- "full analysis, top-down splaying, and link-cut
dynamic trees are Chapter~25."  This chapter pays all three.

A link-cut tree answers a question no static structure can: keep a forest whose
edges you \emph{add and delete online}, and after every change return an aggregate
(sum, min, max) over the unique path between two nodes, or decide whether they are
even connected.  The trick is to store the real tree (the \term{represented tree})
as a handful of splay trees (the \term{auxiliary trees}) -- one per \term{preferred
path} -- and let a single operation, \code{access}, re-shuffle which edges are
preferred so that the root-to-node path you care about is always one contiguous
splay tree you can read in \(O(\log n)\).

\textbf{You will learn to}
\begin{itemize}[nosep]
  \item run the \emph{splay engine} -- splaying by access, the zig / zig-zig /
        zig-zag steps, the potential function that makes one operation \(O(n)\)
        worst-case but \(O(\log n)\) amortised, top-down splaying with no
        recursion, and size-maintained splay trees as an order-statistics set
        (\(\secref{c23-splay}\), solved SPOJ ORDERSET); \hfill\lv{3}--\lv{4}
  \item cut a rooted forest into \emph{preferred paths}, store each as one aux
        splay keyed by depth, walk the \code{access} spine, and read the LCA of
        two nodes with two accesses and no rerooting
        (\(\secref{c23-prefer}\), solved SPOJ DYNALCA); \hfill\lv{4}--\lv{5}
  \item \emph{evert} a tree with \code{makeroot} + a reverse lazy tag, then
        \code{link}, \code{cut}, and fold a vertex monoid up the aux trees to
        answer a path aggregate in \(O(\log n)\)
        (\(\secref{c23-makeroot}\), the worked dynamic path-sum); \hfill\lv{5}
  \item push \emph{beyond paths} -- maintain a \code{vsum} accumulator for virtual
        (non-preferred) children so subtree aggregates survive an access, and see
        why connectivity is often an Euler-tour-tree or rollback-DSU job, not an
        LCT job (\(\secref{c23-beyond}\)); \hfill\lv{5}
  \item read a specification and decide whether the topology is even dynamic --
        if it is fixed, hand the job back to HLD / Euler+Fenwick / centroid
        (\(\secref{c23-boundary}\), the pick table). \hfill\lv{5}
\end{itemize}

\textbf{Prerequisites.} Chapter~8 (\chref{c06}) -- the splay tree and its
amortised behaviour, treated here as the engine we finally drive all the way.
Chapter~16 (\chref{c11}) -- \(\lca\) and the static \(\mathit{tin}\)/depth world a
link-cut tree replaces.  Chapter~22 (\chref{c16}) -- heavy-light decomposition:
the \emph{static} cousin of a preferred-path split, cited (never re-taught) in
\(\secref{c23-prefer}\).  Chapter~5 (\chref{c24}) -- the Euler tour that becomes
an Euler-tour tree in \(\secref{c23-beyond}\).  Chapter~24 (\chref{c18}) -- the
offline divide-and-conquer engine that beats an LCT for a known edge timeline.
\end{chaptertheorybox}

% ---------------- Reading paths + ladder position ----------------
\begin{chaptertheorybox}{How to Read This Chapter (three paths)}
Level badges \(\lv3\)--\(\lv6\): \(\lv3\) is splaying a node to the root of a BST
by hand and counting rotations; \(\lv{4}\) is a size-maintained splay set that
answers k-th / count-less-than, or a correct \code{access} spine on the running
forest; \(\lv{5}\) is \code{makeroot} with a reverse tag driving link/cut/path-sum,
or a subtree \code{vsum} fix-up; \(\lv{6}\) is the Summit -- a block-cut tree with
an LCT path-min riding on top of it.

\textbf{Engine path (\(\lv3\)--\(\lv{5}\), "make splaying second nature").}
\secref{c23-splay} in full -- the three splay cases, the potential, top-down
splaying, size aggregates, solved SPOJ ORDERSET -- then \secref{c23-prefer}, where
the same splay becomes the auxiliary tree of a preferred path (solved SPOJ
DYNALCA), capped by the Summit \(\probref{c23-summit}\).

\textbf{Full-LCT path (\(\lv{4}\)--\(\lv{6}\), "a dynamic forest I can query").}
\secref{c23-prefer} for \code{access} and rooted LCA, then \secref{c23-makeroot}
for evert / link / cut / path aggregate (the worked dynamic path-sum), then
\secref{c23-beyond} for the subtree \code{vsum} that makes a real competitive LCT
correct -- and the Summit, where an LCT sits on a block-cut tree.

\textbf{Boundary path (\(\lv{4}\)--\(\lv{5}\), "is this even a link-cut job?").}
\secref{c23-boundary}: the pick table that hands fixed-tree path ranges to HLD
(\chref{c16}), subtree ranges to Euler+Fenwick (\chref{c24}), offline edge
timelines to rollback-DSU / centroid (\chref{c18}), and pure connectivity to an
Euler-tour tree -- so you only pay the LCT constant when the topology truly moves.
\end{chaptertheorybox}

\begin{figure}[H]
\centering
\begin{tikzpicture}[diagram, x=8.5mm, y=6.5mm]
  \foreach \st/\lab in {1/lv1,2/lv2,3/lv3,4/lv4,5/lv5,6/lv6} {
    \fill[\lab!22, draw=\lab, line width=0.8pt] (\st-1,0) rectangle (\st,\st*0.52);
    \node[font=\scriptsize\bfseries, text=\lab] at (\st-0.5,\st*0.52-0.42) {L\st};}
  \node[annot, anchor=south west, font=\scriptsize]
    at (0.05,0.62) {prereq (Ch.8): a splay tree, rotate-to-root};
  \node[annot, anchor=south west, font=\scriptsize]
    at (1.05,1.14) {25.1 zig / zig-zig / zig-zag; potential; amortised O(log n)};
  \node[annot, anchor=south west, font=\scriptsize]
    at (2.05,1.66) {25.1 top-down splay, size aggregates, solved SPOJ ORDERSET};
  \node[annot, anchor=south west, font=\scriptsize]
    at (3.05,2.18) {25.2 preferred paths, access spine, solved SPOJ DYNALCA};
  \node[annot, anchor=south west, font=\scriptsize, text=lv5]
    at (4.05,2.70) {25.3 makeroot/link/cut/path-sum; 25.4 subtree vsum, ETT};
  \node[annot, anchor=south west, font=\scriptsize, text=lv6]
    at (5.05,3.22) {25.5 the pick table; Summit CF 487E (block-cut tree + LCT)};
\end{tikzpicture}
\caption*{Chapter~25 on the six-step ladder: splaying a node to the root by hand is
\(\lv3\); a size-maintained splay set and a correct access spine are \(\lv{4}\);
\code{makeroot} with a reverse tag and a subtree \code{vsum} are \(\lv{5}\); the
Summit (an LCT path-min on a block-cut tree) is step~6.}
\end{figure}

% ---------------- The running example ----------------
\subsection*{The Running Example: One Rooted Forest, Accessed Over and Over}

Every section in this chapter runs on the same rooted forest \(R\) on \(n=8\)
nodes, fixed once here and authoritative (if a later trace disagrees, this figure
wins).  It is a single tree rooted at \(1\); parent-to-child edges, 1-based:
\[
  1\!\to\!2,\quad 1\!\to\!3,\quad 2\!\to\!4,\quad 2\!\to\!5,\quad
  5\!\to\!6,\quad 5\!\to\!7,\quad 3\!\to\!8 .
\]
Depths from the root are \(\depth(1)=0\), \(\depth(2)=\depth(3)=1\),
\(\depth(4)=\depth(5)=\depth(8)=2\), \(\depth(6)=\depth(7)=3\).  In \(\secref{c23-prefer}\)
we mark the preferred edges (drawn solid) as the ones on the access path to node
\(7\): \(1\!\to\!2\!\to\!5\!\to\!7\) is the preferred path containing the root, and
\(\{4\},\{6\},\{3,8\}\) hang off it as separate preferred paths -- three auxiliary
splay trees plus the root's spine.

\begin{figure}[H]
\centering
\begin{tikzpicture}[diagram, scale=0.95, transform shape,
   tnode/.style={circle, draw=primaryblue!70, fill=primaryblue!8, line width=0.7pt,
                 minimum size=6.5mm, font=\scriptsize\bfseries, inner sep=1pt},
   pref/.style={circle, draw=difforange, fill=difforange!14, line width=1pt,
               minimum size=6.8mm, font=\footnotesize\bfseries, inner sep=1pt}]
  \node[pref] (n1) at (0,0) {1};
  \node[pref] (n2) at (-2.4,-1.4) {2};
  \node[tnode] (n3) at (2.4,-1.4) {3};
  \node[tnode] (n4) at (-3.8,-2.9) {4};
  \node[pref] (n5) at (-1.2,-2.9) {5};
  \node[tnode] (n8) at (2.4,-2.9) {8};
  \node[tnode] (n6) at (-2.2,-4.4) {6};
  \node[pref] (n7) at (-0.2,-4.4) {7};
  \draw[difforange, line width=1.6pt] (n1)--(n2) (n2)--(n5) (n5)--(n7);
  \draw (n1)--(n3) (n2)--(n4) (n3)--(n8) (n5)--(n6);
  \node[annot, font=\scriptsize, text=difforange] at (2.2,0.2) {preferred path to 7 (solid, orange)};
\end{tikzpicture}
\caption{The running forest \(R\) (\(n=8\), rooted at \(1\)).  The orange solid
edges \(1\!-\!2\!-\!5\!-\!7\) are the preferred path after an \code{access(7)}; the
orange nodes are one auxiliary splay tree keyed by depth.  The white subtrees
\(\{4\}\), \(\{6\}\), \(\{3,8\}\) are each their own preferred path (their own aux
tree) whose top node's \emph{path-parent} pointer aims up into the orange spine.}
\label{fig:c23-tree}
\end{figure}

% ---------------- Which job, cost chart ----------------
\subsection*{Which Job? (the question-shape chart)}

\begin{center}
\small
\renewcommand{\arraystretch}{1.25}
\setlength{\tabcolsep}{3pt}
\begin{tabular}{@{}p{0.30\linewidth} p{0.17\linewidth} p{0.29\linewidth} >{\raggedright\arraybackslash}p{0.14\linewidth}@{}}
\toprule
\textbf{Question} & \textbf{Technique} & \textbf{Why not a static structure} & \textbf{Where} \\
\midrule
order statistics on a dynamic set (k-th, count \(<x\)) & size-maintained splay & a sorted array is \(O(n)\) per insert; a BIT needs small keys & \secref{c23-splay} \\
LCA while edges link and cut online & rooted LCT \code{access} & binary lifting is rebuilt in \(O(n\log n)\) per edit & \secref{c23-prefer} \\
aggregate over the path \(u\!-\!v\) as the tree reroots & \code{makeroot}+path fold & HLD assumes a fixed root and fixed chains & \secref{c23-makeroot} \\
\code{link(u,v)}, \code{cut(u,v)}, ``are \(u,v\) connected?'' & LCT / Euler-tour tree & no static index survives a topology edit & \secref{c23-makeroot} \\
subtree aggregate on a \emph{changing} tree & LCT + virtual \code{vsum} & a Euler tour is destroyed by a cut & \secref{c23-beyond} \\
path range query, \emph{fixed} tree & \emph{go to HLD} & the LCT constant is larger; chains are simpler & \chref{c16} \\
offline edge timeline, all queries known & \emph{rollback DSU / centroid} & you never need an online structure & \chref{c18} \\
\bottomrule
\end{tabular}
\end{center}

\begin{insightbox}[One promise decides everything: make the path you need contiguous]
A splay tree is just a BST; a link-cut tree is a \emph{set} of splay trees, each
holding one preferred path of the real forest, ordered by depth.  The only real
operation is \code{access(v)}: splay \(v\) to the top of its path and repeatedly
walk up through path-parent pointers, at each step declaring the edge you came
from \emph{preferred} and the edge you leave behind \emph{unpreferred}.  After one
\code{access}, the whole root-to-\(v\) path is a single aux splay whose in-order
sequence is exactly that path -- so a path aggregate is the splay's subtree fold,
read in \(O(\log n)\).  Because every step is a splay, the amortised bound
inherits the \(O(\log n)\) you proved in Chapter~8, and the tree can be rerooted,
relinked and recut freely -- the represented tree is a \emph{view}, the aux splay
trees are the \emph{data}.
\end{insightbox}

% ---------------- Taxonomy chart ----------------
\begin{figure}[H]
\centering
\begin{tikzpicture}[diagram, scale=0.95, transform shape,
    mnode/.style={circle, draw=#1, line width=0.9pt, fill=#1!8, minimum size=9mm,
                  font=\scriptsize\bfseries, text=#1!60!black}]
  \node[mnode=primaryblue] (a) at (0,0)   {splay};
  \node[annot, font=\scriptsize, anchor=west, text width=9.8cm, align=left] at (1.7,0) {rotate the touched node to the root; zig / zig-zig / zig-zag; the potential makes a run O(log n) amortised; keep subtree sizes for order statistics (25.1, SPOJ ORDERSET)};
  \node[mnode=successgreen] (b) at (0,-1.1) {access};
  \node[annot, font=\scriptsize, anchor=west, text width=9.8cm, align=left] at (1.7,-1.1) {split the rooted forest into preferred paths, one aux splay each; access(v) makes root-to-v one contiguous splay; two accesses read the LCA, no reroot (25.2, SPOJ DYNALCA)};
  \node[mnode=difforange]   (c) at (0,-2.2) {evert};
  \node[annot, font=\scriptsize, anchor=west, text width=9.8cm, align=left] at (1.7,-2.2) {makeroot(v) = access(v) + flip a reverse lazy tag on v's splay; now v is the tree root, so link(u,v), cut(u,v) and a path fold u--v all become access-and-read (25.3, worked path-sum)};
  \node[mnode=warningred]   (d) at (0,-3.3) {virtual};
  \node[annot, font=\scriptsize, anchor=west, text width=9.8cm, align=left] at (1.7,-3.3) {a splay child is only the PREFERRED child; a subtree aggregate needs a separate vsum over the non-preferred (virtual) children, fixed up every time access detaches a solid edge (25.4); connectivity -> Euler-tour tree / DSU};
  \node[mnode=softpurple]   (e) at (0,-4.4) {pick};
  \node[annot, font=\scriptsize, anchor=west, text width=9.9cm, align=left] at (1.7,-4.4) {is the topology truly dynamic? yes -> LCT; fixed path ranges -> HLD; fixed subtree ranges -> Euler+Fenwick; offline timeline -> rollback DSU/centroid; only connectivity -> ETT (25.5)};
\end{tikzpicture}
\caption{The chapter's five-state taxonomy.  \emph{splay} is the engine; \emph{access}
turns the engine into a rooted dynamic-LCA structure; \emph{evert} adds rerooting
and so full link/cut/path-aggregate power; \emph{virtual} is the extra accumulator a
real LCT needs for subtree queries and the connectivity-vs-path split; \emph{pick}
is the boundary choosing between an LCT and its cheaper neighbours.  Each solved
problem sits on one row: ORDERSET on \emph{splay}, DYNALCA on \emph{access}, the
worked path-sum on \emph{evert}, the \code{vsum} example on \emph{virtual}, and the
Summit CF~487E on \emph{evert}+a block-cut condensation.}
\label{fig:c23-taxonomy}
\end{figure}

% ---------------- Who owns what ----------------
\subsection*{Who Owns What (so nothing is learned twice)}

\begin{center}
\small
\renewcommand{\arraystretch}{1.25}
\begin{tabular}{@{}p{0.30\linewidth} p{0.24\linewidth} p{0.40\linewidth}@{}}
\toprule
\textbf{Idea} & \textbf{Owner} & \textbf{This chapter's share} \\
\midrule
splay tree definition, rotate-to-root & Ch.~8 (\chref{c06}) & the IOU is paid: full potential analysis, top-down splaying and size aggregates (25.1); the single rotation is never re-drawn \\
static heavy-light path split & Ch.~22 (\chref{c16}) & the preferred path is the SAME heavy-path idea made dynamic; cited once in 25.2, boundary drawn in 25.5 \\
\(\lca\) by binary lifting & Ch.~16 (\chref{c11}) & 25.2's two-access LCA is an ALTERNATIVE for a moving forest, not a re-derivation \\
Euler tour, flatten to array & Ch.~5 (\chref{c24}) & becomes an Euler-tour \emph{tree} for connectivity in 25.4; tin/tout never re-taught \\
offline divide-and-conquer on a timeline & Ch.~24 (\chref{c18}) & 25.5 routes a known edge timeline to rollback-DSU / centroid instead of a live LCT \\
dynamic connectivity (only) & Ch.~17 DSU (\chref{c28}) & if you never need a path value, a (rollback) DSU beats an LCT; forward-cited in 25.4 \\
compressed tree of a marked node set & Ch.~26 (\chref{c30}) & explicitly out; virtual trees are the next chapter \\
\bottomrule
\end{tabular}
\end{center}

% ---------------- Recognition guide ----------------
\begin{recognitionbox}[Pattern card -- "the tree itself changes under my queries"]
\reckey{statements with an online \emph{edge edit} on a forest: ``\code{link(u,v)}
/ \code{cut(u,v)} / query the path sum from \(u\) to \(v\), all interleaved'',
``support rerooting the tree, then a path aggregate'', ``after adding or deleting a
road, is \(u\) still reachable from \(v\)? and what is the cheapest city on the
route?'', ``maintain a dynamic forest and answer a value that lives on the unique
path between two live nodes''.  The tell is a \emph{topology edit} (an edge appears
or vanishes) interleaved with a \emph{path} query -- if the edges are fixed, you
wanted HLD; if you only need connectivity, you wanted DSU.}
\recuse{
  \item edges link and cut online \emph{and} you query a value along the current path -- full LCT with \code{makeroot} (\(\secref{c23-makeroot}\));
  \item a rooted forest where only links/cuts change and you need LCA -- rooted LCT, no evert (\(\secref{c23-prefer}\));
  \item a dynamic ordered set / sequence needing k-th or rank -- size-maintained splay, the engine alone (\(\secref{c23-splay}\)).
}
\recnotuse{
  \item the tree is \emph{fixed} and you query/update a path range -- HLD (\chref{c16}), smaller constant;
  \item the tree is \emph{fixed} and you query a subtree -- Euler tour + Fenwick (\chref{c24});
  \item all edge edits are \emph{known in advance} (offline) -- rollback DSU or segment-tree-on-time, no online structure;
  \item you only need \emph{connectivity}, never a path value -- (Euler-tour) DSU is simpler (\(\secref{c23-beyond}\));
  \item a \emph{subset of marked nodes} compressed to its Steiner tree -- virtual trees (\chref{c30}).
}
One-line test: \emph{an edge appears or disappears between two queries AND the
answer depends on the current path -- then a link-cut tree; otherwise a static
engine is faster.}
\end{recognitionbox}

% ---------------- Constraint -> technique ----------------
\subsection*{Reading the Constraints (which technique the limits want)}

\begin{center}
\small
\renewcommand{\arraystretch}{1.25}
\begin{tabular}{@{}p{0.34\linewidth} p{0.40\linewidth} p{0.20\linewidth}@{}}
\toprule
\textbf{Constraint signature} & \textbf{What it rules in/out} & \textbf{Technique} \\
\midrule
dynamic set, keys \(|x|\le10^9\), k-th / rank queries & a BIT cannot index huge keys; a splay holds them and carries sizes & \secref{c23-splay} \\
link/cut a forest, ask LCA, \(N,M\le10^5\) & binary lifting rebuilds per edit; \code{access} answers in \(O(\log n)\) & \secref{c23-prefer} \\
\code{link}, \code{cut}, path sum/min interleaved online & only a rerootable aux-splay forest survives; LCT & \secref{c23-makeroot} \\
subtree sum on a tree that also link/cuts & preferred children miss the off-path subtrees; add a virtual \code{vsum} & \secref{c23-beyond} \\
edge insert/delete timeline given \emph{offline}, connectivity only & no online structure needed; rollback DSU / D\&C (\chref{c18}) & \secref{c23-boundary} \\
\bottomrule
\end{tabular}
\end{center}

% ---------------- Notation ----------------
\subsection*{Notation Used Only Here}

The represented forest is \(R=(V,E)\), \(|V|=n\), each tree rooted (until an
\code{evert} changes which node is root).  A node's \emph{aux} children are the
at-most-two children in its splay tree (the \emph{preferred} path neighbours, one
shallower and one deeper); its \emph{path-parent} is a single pointer from the top
of a preferred path to the node just above it in the represented tree.  A node
\(x\) is an \term{aux root} (\code{is\_root(x)}) when it is not a child in any
splay tree -- equivalently its parent, if any, treats it as a path-parent, not an
aux child.  \code{access(v)} detaches \(v\)'s right (deeper) preferred child and
repeatedly climbs path-parents, so after it \(v\)'s splay holds the whole current
root-to-\(v\) path in depth order.  \code{makeroot(v)} = \code{access(v)} then a
reverse lazy tag on \(v\)'s splay, making \(v\) the represented root.  The
aggregate of an aux splay node \(x\) is \(\mathit{sum}(x)\) (or \(\min\)) over its
splay subtree -- i.e.\ over the preferred path segment \(x\) represents; a
\emph{virtual child} is a represented child not currently on \(x\)'s preferred
path, and their values are held in \(\mathit{vsum}(x)\) (\(\secref{c23-beyond}\)).
\(\depth(\cdot)\) is depth in the current represented tree.  Every section runs on
\(R\) above; if a table here and a trace there disagree, this opener is
authoritative.
